\section*{Chapter \ref{ch:s1:anatomy-of-a-stat-arb-book} --- Anatomy of a Stat-Arb Book}

\begin{solution}{exo:s1:anatomy-of-a-stat-arb-book:1}
Independent bets add in variance: $\sqrt{946 \times 252} \times 0.02 = 9.8$ positions' worth, against $946 \times 252 = 238\,392$ position-days. The
edge adds linearly and the noise with the square root, which is the whole case for breadth.
\end{solution}

\begin{solution}{exo:s1:anatomy-of-a-stat-arb-book:2}
$0.04 - 0.0342 = 0.0058$: 0.58\% of capital a year, the price of financing longs beyond capital and of the fees inside the rebate.
\end{solution}

\begin{solution}{exo:s1:anatomy-of-a-stat-arb-book:3}
$0.5 \times 1.5 + 0.5 \times 1.5 = 1.5$ times capital, which the book does not have. Equity must be half the gross, so the gross is at most $1/0.5 = 2.0$ times
capital; a requirement of 15\% of the gross allows $1/0.15 = 6.67$ times.
\end{solution}

\begin{solution}{exo:s1:anatomy-of-a-stat-arb-book:4}
The rebate is $1.5 \times (0.04 - 0.0025)$, the long financing $0.5 \times 0.045$: $1.5 \times 0.0375 - 0.5 \times 0.045 = 0.0337$, 3.37\% of capital a
year. The chapter's book earns 3.42\% because its capped longs sum to slightly less than 1.5 on some days.
\end{solution}

\begin{solution}{exo:s1:anatomy-of-a-stat-arb-book:5}
The signal's exposure to the styles changes every day (momentum drifts, sizes and prices move), and each day's neutralisation removes it again by trading. The
book neutral only to beta and industries simply carries the style exposure and does not trade it.
\end{solution}

\begin{solution}{exo:s1:anatomy-of-a-stat-arb-book:6}
If $f_t$ is the cross-sectional regression's coefficient and $e_t = r_t - X f_t$ its residual, then $r_t = X f_t + e_t$ by construction and
$w^\top r_t = w^\top X f_t + w^\top e_t = (X^\top w)^\top f_t + w^\top e_t$. It breaks when the exposures used in the attribution differ from those of the
regression (a different day's, or a different model's), when names missing from the regression are held, or when the book's return is measured on different
prices from the model's.
\end{solution}

\begin{solution}{exo:s1:anatomy-of-a-stat-arb-book:7}
With a half-life of one day the book turns over 129\% of capital a day, pays 32.5\% a year in costs and returns 27.5\% at a Sharpe ratio of 6.16; with five days
56\%, 14.0\%, 11.7\% and 2.47. Here faster is better, because the simulated reversal is strong and short-lived and the cost of five basis points is low: slowing
the book loses more signal than it saves in costs. At realistic costs the ranking can reverse (chapter~2).
\end{solution}

\begin{solution}{exo:s1:anatomy-of-a-stat-arb-book:8}
Zero beta and zero dollar exposure leave every other factor free. The chapter's book neutral to beta and industries had 99.7\% of its variance from factors,
nearly all momentum; its return was a style bet. Market neutral is not factor neutral, and alpha is what remains after all the factors the firm can buy
cheaply.
\end{solution}

\begin{solution}{pb:s1:anatomy-of-a-stat-arb-book:1}
\begin{enumerate}
\item An equal blend of the z-scored one-day residual reversal and 12--1 month momentum, smoothed (half-life two days), neutralised by regression on the risk
model's exposures, scaled to gross three, capped at 1\% a name, with no shorts in hard-to-borrow names.
\item 452 long and 495 short on average; the median position is 0.25\% of capital, the largest 1.0\%.
\item The information ratio grows with the square root of the number of independent bets; the law assumes independent bets and ignores costs, and the book's
common exposures and daily trading violate both.
\item 1.9 million.
\item 23.6\%, 23.6\% and 1.00 neutral to beta and industries; 22.5\%, 4.5\% and 5.03 also neutral to the styles.
\item 99.7\% in the first, almost all momentum (a momentum P\&L of 9.0\% a year with a volatility of 23.1\%, and value $-2.9\%$); 7.6\% in the second.
\item Turnover of 91\% of capital a day against 42\%, and trading costs of 22.9\% a year against 10.6\%.
\item The simulation plants a strong reversal and a persistent drift at costs of five basis points; published stat-arb Sharpe ratios after costs are nearer 1.
\item Rebate 5.6\%, long financing $-2.2\%$, borrow fees $-0.4\%$ (inside the rebate): net $+3.4\%$ of capital a year.
\item Equity of half the gross, so at most twice capital; a gross of three needs risk-based portfolio margin (a 15\% stress without netting requires 0.45 of
capital).
\item They are flagged before optimisation and not shorted; had they been, their fees (median around 3\%) would have cut the rebate and exposed the book to recall.
\item \textbf{Named result.} 99.7\% of the P\&L variance of the book not neutral to the styles came from factor exposures; the financing drag (long financing and
borrow fees) was 6.2\% of the neutral book's gross alpha of 42.0\%.
\item Sharpe ratios after costs of 1.44 for 1997--2007 for their PCA strategies, 0.9 in 2003--2007; the ETF version degraded similarly after 2002.
\item Books built this way held similar positions and lost together when some of them were unwound.
\item Data checked and stored, risk model updated, signals blended and smoothed, the book optimised, trades executed the next session, financing accrued, the P\&L
attributed and reconciled.
\item Because a part that does not add up means a wrong position, price, exposure or charge: an unexplained residual is a defect, not a line in the report.
\item Realistic costs with impact (Book~7, chapter~27), locates and recalls, a real risk model, and signals whose strength is estimated, not planted.
\item Attribute its P\&L: a large share of variance from factor returns, or a high correlation with style indices, means a factor bet.
\item What share of its variance comes from factors it does not claim to hold.
\item A long--short portfolio of many small, individually unreliable bets, neutral to what it does not want to hold, whose return comes from their number.
\end{enumerate}
\end{solution}

\begin{solution}{iq:s1:anatomy-of-a-stat-arb-book:1}
Zero net dollars and zero beta, so the market's move does not move it. It is still exposed to industries, styles (momentum, value, size), crowding with similar
books, and financing and borrow conditions.
\end{solution}

\begin{solution}{iq:s1:anatomy-of-a-stat-arb-book:2}
Longs with capital and a debit balance at the prime broker; shorts with borrowed stock, the proceeds held as collateral earning the benchmark rate less the fee.
What goes wrong: fees rise or names are recalled, margin requirements rise in a stress, the broker cuts financing, and the book is forced to shrink at the worst
time.
\end{solution}

\begin{solution}{iq:s1:anatomy-of-a-stat-arb-book:3}
The information ratio is approximately the information coefficient times the square root of breadth. A small forecasting skill applied to many independent
names and rebalanced often gives a high ratio, which is stat arb's premise; its limits are the bets' correlation and the costs of trading them.
\end{solution}

\begin{solution}{iq:s1:anatomy-of-a-stat-arb-book:4}
More capital in the same signals (the edge is shared and the costs rise), tighter spreads and faster markets that removed some mispricings, and crowding that
made the books similar enough to lose together.
\end{solution}

\begin{solution}{iq:s1:anatomy-of-a-stat-arb-book:5}
Point-in-time positions, prices, the risk model's exposures and factor returns; P\&L split into factor (by factor), specific, costs (by trade) and financing (by
name). It must reconcile with positions times returns and with the prime broker's statement of cash, fees and rebates.
\end{solution}

\begin{solution}{iq:s1:anatomy-of-a-stat-arb-book:6}
Each position is $G/N$; the specific variance is $N (G/N)^2 \sigma^2$, so the volatility is $G\sigma/\sqrt N$. The expected return is $N (G/N)\alpha = G\alpha$,
and the ratio is $\alpha\sqrt N/\sigma$: it grows with $\sqrt N$ at a fixed per-name edge, independent of $G$.
\end{solution}
